\chapter{Conclusion and Outlook}\label{chap:conclusion}

This thesis investigated whether point-cloud upsampling improves downstream
perception, using object classification on ModelNet40 and three-dimensional
object detection on KITTI. The two-line design separated densification of
an available observation from recovery after fourfold sparsification.
Five upsampling methods completed the HPC object-level study. The lab
scene-level study combined earlier frozen-detector comparisons with a
later, fully evaluated PU-GCN input and detector-adaptation experiment.
The combined evidence shows that added coordinates, geometric agreement,
and useful task information are distinct properties.

\section{Findings from the Completed Experiments}

On ModelNet40, none of the five upsampled Line~A representations exceeded
the original 1,024-point classifier. Original best reported OA was 91.95\%, while
the highest upsampled best OA was 91.63\% for PU-GCN. After sparsification
to 256 points, PU-Net alone improved the selected-checkpoint result,
from 90.85\% to 91.27\%. Its final-epoch gain was smaller, from 90.46\%
to 90.65\%. These are results from separately trained, single-seed
classifiers, not repeated-run estimates
(Tables~\ref{tab:mn-a} and~\ref{tab:mn-b}).

The geometric ranking did not reproduce the classification ranking.
PU-GCN had the lowest mean equal-cardinality Chamfer distance in both
lines, while PU-Net had the best Line~B classification and Hausdorff
distance. PointNet++ reception measurements further distinguished
stored density from the support used by local feature aggregation.
PU-GCN's CD advantage held for more than 93\% of test objects in
each line, while PU-Net had the lowest Line~B HD for about 90\%.
Thus the geometric and task mismatch was not explained only by a few
outliers. The two retained point-cloud examples illustrate that
distinction without assigning unrecorded classification outcomes
(Sections~\ref{sec:controls}, \ref{sec:equaln-geometry}
and~\ref{sec:modelnet-visual}).

The KITTI study first evaluated 20 three-epoch or frozen-weight arms
on all 3,769 validation frames, then extended training for the four
observed-first detector--line combinations. The extended schedules
improved on both frozen observed-first inputs and the earlier adapted
checkpoints, while satisfying the specified Car AP plateau criterion.
Selected Line~A scores reached 71.59 for PointRCNN and 75.97 for
CenterPoint, below the original baselines of 81.95 and 79.28.
Line~B reached 57.05 and 62.90, below the adapted sparse baselines of
68.33 and 68.05 (Table~\ref{tab:k-extended-comparison}).
Thus the residual deficit persisted beyond the original short
adaptation schedule. PU-GCN itself remained fixed throughout.

The integration diagnostics explain why a method name alone is insufficient
to interpret a scene-level result. Local patch construction, correct
normalisation, measurement preservation, and the detector's point or voxel
budget all changed the evaluated input. Correcting PU-Net normalisation
and patch locality substantially improved its diagnostic result. Enlarging
PointRCNN input or changing CenterPoint occupancy did not yield a universal
remedy. These findings support evaluation of the complete implemented
pipeline, with the scope of every control stated explicitly
(Sections~\ref{sec:patch-diagnostics} and~\ref{sec:budget-results}).

\section{Limits of the Evidence}

The ModelNet40 classifiers used one seed and selected their best checkpoint
on the test set. Their saved OA and class scores average batch-level
ratios; they are not exact object-count proportions. Initial mesh sampling
also used Python's process-dependent hash in its seed. Reusing the archived
point sets supports repeatability of the completed runs, whereas exact
regeneration from meshes requires additional hash-state control. Reporting epoch~200 alongside best OA reduces dependence
on that selection in the presentation, but does not create an independent
validation protocol. Four reused training objects limit the mesh-reference
classifier controls. Equal-cardinality mesh-reference comparisons also
use independently normalised point samples rather than one shared
mesh-derived transform. Their distance rankings are conditional on that
coordinate convention. P2F used independent cloud and mesh normalisation,
and NUC query centres were not perfectly shared across methods; these
remain secondary geometry diagnostics. Missing per-object classifier
predictions prevent paired prediction-level tests
(Sections~\ref{sec:fc-4-3-1}, \ref{sec:fc-4-10-1},
and~\ref{sec:fc-4-10-2}).

The KITTI extension established a plateau under a 0.2-AP, three-epoch
patience rule for Car Moderate 3D AP, rather than global optimisation
convergence or convergence of all classes. Checkpoints were selected and
reported on the same validation set. Neither the original nor the sparse
baseline was extended to an equal training budget; Line~A also has no
adapted original baseline. These design limits prevent attributing the
system-level differences to generated geometry alone. The earlier
PU-Net full-validation pipeline had a normalisation error, and the lab
PU-EdgeFormer execution used a compatibility path. Historical evaluator
variants and small diagnostic subsets consequently cannot be combined
with the final matrix as one architecture ranking. Development also
examined validation outcomes; these results are not a new blind test
(Sections~\ref{sec:adaptation-evidence} and~\ref{sec:kitti-diagnostics}).

\section{Implications and Future Work}

The practical implication is to assess upsampling after the actual task
model and its input adapter. A useful evaluation records visible
measurements, generated candidates, strict output, and effective model
input separately. It also distinguishes geometry references and training
regimes. Such reporting makes a limited recovery gain interpretable
without promoting it into a general claim about point-cloud densification.

Further experiments could strengthen the conclusions by separating
classifier validation from the fixed test set, repeating training across
seeds, regenerating all mesh-reference training controls, evaluating
geometry with a shared mesh-derived transform, and retaining
per-object predictions. For KITTI, the remaining comparison should use
equally trained original, sparse and upsampled inputs, together with
independent checkpoint selection and final evaluation. Repeated seeds
and alternative optimisation schedules would test the stability of the
residual deficit beyond the completed plateau assessment. Native implementations
under corrected common patch processing would permit cleaner method
comparisons. These are proposed extensions, not experiments performed in
this thesis.

Within the executed scope, geometric improvement was insufficient to
establish downstream improvement. PU-Net provided a limited object-level
recovery benefit, while observation preservation and detector adaptation
partially recovered KITTI performance. The evidence supports these
conditional findings rather than a universal benefit or universal failure
of point-cloud upsampling.
